%*###############################################################
\section{Power and polarization stability results}
\label{sec:power_polarization_results}
%*###############################################################

The black-fiber injection system described in Chapter~2 provides the
counterpropagating optical field required to generate the conveyor-belt
potential. For this application, stable optical power and polarization are
essential. Power fluctuations directly affect the lattice depth and may lead to parametric heating, whereas
polarization fluctuations may reduce the interference contrast and,
consequently, the reproducibility of the optical potential.

This section presents the experimental characterization of the power and
polarization stability of the optical field transmitted through the black
fiber. The system is analyzed both in the time domain, using the fully
overlapping Allan deviation, and in the frequency domain, using the
Amplitude Spectral Density (ASD). These complementary analyses make it
possible to identify the relevant noise processes, determine the timescales
over which the system is most stable, and compare the measured noise with
the theoretical shot-noise limit.

%*###############################################################
\subsection{Time-domain stability}
%*###############################################################

The time-domain stability of the transmitted optical power and polarization
was evaluated from a continuous acquisition lasting $600~\mathrm{s}$. 
Figure~\ref{fig:time_series_raw} displays the raw voltage signals for the power channels and the corresponding fluctuation of the polarization angle recorded during this time window. 
The fully overlapping Allan deviation was then calculated for the recorded signals
in order to distinguish short-term fluctuations from slow environmental
drifts.

\begin{figure}[htbp]
    \centering
    \includegraphics[width=0.8\textwidth]{Cap_3_figures/data_SAMPLE_600s_AfterBlackFiber_time_series.png}
    \caption{Time series of the raw signals measured after the black fiber during a $600~\mathrm{s}$ acquisition. Top: voltage amplitude of the two transmission channels (CH1 and CH2). Bottom: corresponding fluctuation of the polarization angle $\Delta\theta$.}
    \label{fig:time_series_raw}
\end{figure}

The top panel of Figure~\ref{fig:allan_deviation_results} shows the
fractional Allan deviation $\sigma_y(\tau)$ of the optical power. The total
power, obtained as the sum of the transmitted and reflected channels,
reaches a minimum fractional Allan deviation of approximately
$8\times10^{-4}$ at an averaging time of $\tau \approx 0.1~\mathrm{s}$.
For shorter averaging times, $\sigma_y(\tau)$ follows the characteristic $1/\sqrt{\tau}$ dependence associated with white noise. For
$\tau > 0.1~\mathrm{s}$, the Allan deviation increases progressively,
exceeding $10^{-2}$ at $\tau = 200~\mathrm{s}$. This behavior indicates
that the long-term power stability is limited by low-frequency processes,
most likely associated with thermal drifts and slow mechanical relaxation
of the optical components.

The bottom panel of Figure~\ref{fig:allan_deviation_results} shows the
angular Allan deviation $\sigma_\theta(\tau)$ of the polarization angle.
The minimum value, approximately $7\times10^{-5}~\mathrm{rad}$, is obtained
at an averaging time of $\tau \approx 3~\mathrm{s}$. At shorter averaging
times, the Allan deviation follows the characteristic
$1/\sqrt{\tau}$ dependence associated with white noise. For
$\tau > 3~\mathrm{s}$, slow environmental drifts become dominant and
$\sigma_\theta(\tau)$ increases, reaching approximately
$8\times10^{-4}~\mathrm{rad}$ at $\tau = 200~\mathrm{s}$.

\begin{figure}[htbp]
    \centering
    \includegraphics[width=0.8\textwidth]
    {Cap_3_figures/data_SAMPLE_600s_AfterBlackFiber_allan_dev.png}
    \caption{Fully overlapping Allan deviation of the power and polarization
    signals measured after the black fiber over a $600~\mathrm{s}$
    acquisition. Top: fractional Allan deviation $\sigma_y(\tau)$ of the
    individual power channels and their sum. The total power reaches a
    minimum fractional Allan deviation of approximately $8\times10^{-4}$ at
    $\tau \approx 0.1~\mathrm{s}$. Bottom: angular Allan deviation
    $\sigma_\theta(\tau)$ of the polarization angle, with a minimum of
    approximately $7\times10^{-5}~\mathrm{rad}$ at
    $\tau \approx 3~\mathrm{s}$. At longer averaging times, both quantities
    are increasingly affected by slow environmental drifts.}
    \label{fig:allan_deviation_results}
\end{figure}

%*###############################################################
\subsection{Frequency-domain analysis}
%*###############################################################

The Allan-deviation analysis identifies the timescales over which the
transmitted field is most stable, but it does not directly reveal the
spectral distribution of the underlying fluctuations. The acquired signals
were therefore also analyzed in the frequency domain by calculating the
Amplitude Spectral Density (ASD) of the photocurrent. The measured spectra
were compared with the theoretical shot-noise levels calculated using the
photocurrents recorded in the corresponding detection channels.

Figure~\ref{fig:fft_low_freq} shows the low-frequency spectra obtained from
the $600~\mathrm{s}$ acquisition over the range from $10^{-3}$ to
$10^2~\mathrm{Hz}$. Below approximately $1~\mathrm{Hz}$, the spectra are
dominated by excess technical noise and exhibit a clear low-frequency
increase consistent with the long-term drifts observed in the Allan
deviation. As the frequency increases, the measured noise decreases, although it remains significantly above the calculated white-noise floors, represented by the dashed lines, clearly showing the presence of technical noise.

\begin{figure}[htbp]
    \centering
    \includegraphics[width=0.8\textwidth]
    {Cap_3_figures/data_SAMPLE_600s_AfterBlackFiber_fft_CH1.png}
    \caption{Amplitude Spectral Density of the photocurrent fluctuations
    measured after the black fiber over the frequency range from
    $10^{-3}$ to $10^2~\mathrm{Hz}$. The dashed lines represent the
    calculated white-noise floors for the corresponding detection channels.
    Excess technical noise dominates the spectra below approximately
    $1~\mathrm{Hz}$.}
    \label{fig:fft_low_freq}
\end{figure}

An additional measurement based on a shorter $600~\mathrm{ms}$ acquisition with a larger bandwidth was performed to characterize the system at frequencies up to $10^5~\mathrm{Hz}$. The
resulting transmitted-photocurrent spectrum is shown in
Figure~\ref{fig:fft_high_freq}. Residual technical noise is visible below
approximately $10^3~\mathrm{Hz}$. Above this frequency, the spectrum flattens but remains well above the calculated shot-noise level of $10^{-11}~\mathrm{A/\sqrt{Hz}}$. This result shows that a residual broadband technical noise is still present at higher frequencies, and the system is not yet limited by the fundamental detection noise, even though the most prominent excess noise is confined to the lower-frequency region.

\begin{figure}[htbp]
    \centering
    \includegraphics[width=0.8\textwidth]
    {Cap_3_figures/data_SAMPLE_600ms_AfterBlackFiber_fft_CH1.png}
    \caption{Amplitude Spectral Density of the transmitted photocurrent up
    to $10^5~\mathrm{Hz}$. The red dashed line represents the calculated
    photocurrent shot-noise level of
    $10^{-11}~\mathrm{A/\sqrt{Hz}}$. Above approximately
    $10^3~\mathrm{Hz}$, the measured spectrum flattens but remains above the shot-noise
    floor.}
    \label{fig:fft_high_freq}
\end{figure}

The time- and frequency-domain measurements provide a consistent
characterization of the stability of the optical field delivered through
the black fiber. FARE CONCLUSiONI PIU QUANTITATIVE CON CONFRONTO